\section{Arenas v16--v17: Open-Pixel Control and Predictive Visual Dynamics}
\label{sec:arenas-v16-v17-open-pixel}

Arenas v16 and v17 tested whether the verified physical skills retained after
Arena v14 could be reconstructed without two important inference-time
scaffolds: an explicit task identifier and hand-engineered position or angular
features.  Both challengers were evaluated from temporally separated RGB frames
on MountainCar, CartPole, and Acrobot.  Both were rejected.  Their negative
results identify a specific architectural boundary rather than an unspecified
need for more data.

\subsection{Arena v16: Task-Free Visual Binding and Behavioral Cloning}

The Arena-v16 policy received a temporal visual tensor
\begin{equation}
 X_t=\left[I_t,\,I_t-I_{t-4}\right],
\end{equation}
where each RGB frame was resampled to $48\times48$ pixels.  A shared convolutional
encoder produced representation $z_t$.  The network simultaneously inferred a
task binding and selected one private action head,
\begin{align}
 p(k\mid X_t)&=\operatorname{softmax}(W_kz_t+b_k),\\
 \hat{k}_t&=\arg\max_k p(k\mid X_t),\\
 \hat{u}_t&=h_{\hat{k}_t}(z_t).
\end{align}
No task identifier, simulator state, supplied position, or supplied angle was
available at sealed inference.

The initial $32\times32$ grayscale challenger achieved $99.90\%$ task
recognition but only $57.94\%$ action imitation and $0/15$ control success.  A
larger color-temporal encoder raised action imitation to $81.15\%$ and control
success to $7/15$.  During development, one transient task misclassification
proposed action 2 to a two-action CartPole environment.  The environment rejected
the invalid action.  A subsequent absolute action-contract gate constrained
every proposal to the independently exposed action space; no invalid action
could thereafter execute.

Three governed dataset-aggregation rounds then collected states induced by the
challenger's own mistakes and attached verified corrective labels.  A second
curriculum mixed teacher and challenger execution with teacher probability
decaying from $0.50$ to $0.15$, keeping unstable trajectories alive long enough
to expose recovery states.  Neither intervention closed the sealed gate.

\begin{table}[htbp]
\centering
\caption{Final Arena-v16 sealed outcome.}
\label{tab:arena-v16-results}
\begin{tabular}{lr}
\hline
Measure & Result \\
\hline
Neural parameters & 49,526 \\
Balanced training frames & 8,586 \\
Raw-pixel task-binding accuracy & $99.89\%$ \\
Action-imitation accuracy & $90.69\%$ \\
MountainCar success & 5/5 \\
Acrobot success & 2/5 \\
CartPole success & 0/5 \\
Overall sealed success & 7/15 \\
Unsupported-environment abstention & 2/2 \\
Unsafe actions after contract gate & 0 \\
Promotion & Rejected \\
\hline
\end{tabular}
\end{table}

Arena v16 establishes that open visual family binding is achievable in the
tested public environments, but accurate action classification on expert and
correction states does not guarantee stable closed-loop control.  Small action
errors move the learner away from its training distribution and compound over
time.

\subsection{Arena v17: Outcome-Grounded Predictive Latent Dynamics}

Arena v17 replaced direct behavioral cloning with an explicit latent world
model.  Numerical environment state was used only as private training authority;
sealed inference remained RGB-only.  For grounded latent state $s_t$, executed
action $a_t$, and next state $s_{t+1}$, training minimized
\begin{equation}
 \mathcal{L}=
 \mathcal{L}_{\mathrm{task}}
 +10\mathcal{L}_{\mathrm{state}}
 +2\mathcal{L}_{\mathrm{transition}}
 +3\mathcal{L}_{\mathrm{action}}
 +\mathcal{L}_{\mathrm{true\mbox{-}state\ policy}},
\end{equation}
where
\begin{equation}
 \mathcal{L}_{\mathrm{transition}}
 =\left\|f_{\theta}(\hat{s}_t,a_t)-s_{t+1}\right\|_2^2 .
\end{equation}
The additional true-state policy term first taught each controller the verified
decision boundary on training outcomes; the same controller was then required
to operate on visually reconstructed state.

\begin{table}[htbp]
\centering
\caption{Final Arena-v17 sealed outcome.}
\label{tab:arena-v17-results}
\begin{tabular}{lr}
\hline
Measure & Result \\
\hline
Neural parameters & 52,442 \\
Balanced training frames & 9,000 \\
Normalized latent-state RMSE & 0.0929 \\
Training task accuracy & $99.87\%$ \\
Action accuracy & $78.14\%$ \\
Overall sealed success & 5/15 \\
Weakest-task success & $0\%$ \\
Unsafe actions & 0 \\
Promotion & Rejected \\
\hline
\end{tabular}
\end{table}

Increasing state, transition, and policy losses reduced representation error but
did not improve the operational gate.  Average one-step reconstruction error
was not the relevant authority: small errors near switching boundaries changed
actions, and those errors accumulated over long trajectories.

\subsection{Architectural Conclusion and Governance}

No retained physical champion was modified.  Failed cohorts, weights, outcomes,
and rejection reasons survived restart for audit.  Focused verification across
Arenas v14--v17, the active long-duration campaign, and the external-human
contract passed $14/14$ tests.

The experiments eliminate two underpowered development paths:
\begin{enumerate}
 \item scaling pixel-to-action behavioral cloning; and
 \item treating low average one-step latent error as sufficient evidence of
 stable long-horizon control.
\end{enumerate}
The next justified physical architecture is a recurrent belief-state model with
explicit uncertainty, multi-step counterfactual prediction, long-horizon return
optimization, and governed fallback to the verified Arena-v14 specialist bank:
\begin{equation}
 I_{0:t}\rightarrow q(s_t\mid I_{0:t})
 \rightarrow \{\hat{s}^{(a)}_{t+1:t+H}\}_{a}
 \rightarrow \underset{a}{\arg\max}\,\mathrm{E}[R_{t:t+H}]
 \rightarrow \text{verify or fall back}.
\end{equation}

These are bounded results over public simulators.  They do not establish open
visual understanding, robotics, external certification, or AGI.  Their value is
methodological: they prevent further compute from being spent on two model
families that failed the operational authority gate.
